% GENERATED FILE. Edit canonical lessons, then run npm run generate:books.
% canonical-source-hash: fnv1a64:e461f454a1179011
% canonical-lessons: BN-C69-goru, BN-C69-chhagol, BN-C69-bhera, BN-C69-pakhi, BN-C69-idur

\chapter{Cow, Goat, Sheep, Bird, Mouse}
\label{ch:bn-cow-goat-sheep-bird-mouse}
\hlchaptermodality{69}

\begin{chapteropening}
\textbf{By the end of this chapter:} \emph{I can say a cow, a goat, a sheep, a bird and a mouse.}

Complete the last lesson of chapter 69: I can say a cow, a goat, a sheep, a bird and a mouse.
\end{chapteropening}

\section[\texorpdfstring{\bn{গরু} (gôru)}{গরু (gôru)}]{\bn{গরু} (gôru) — a cow}

\label{lesson:BN-C69-goru}



\begin{quote}
Before the new one: say the Bengali for a dog, then the Bengali for a cat.
\end{quote}

\subsection*{You'll want to know: \bn{গরু}}
\textbf{\bn{গরু}} — \emph{gôru} — \textquotedblleft{}a cow\textquotedblright{}.

\begin{cousinweb}
\textbf{\bn{গরু}} (\emph{gôru}) is a tadbhava, a word inherited from Sanskrit and reshaped by sound change. It comes from \emph{gorūpa}, \textquotedblleft{}cow-shaped\textquotedblright{} (CDIAL 4313). It is \emph{go}, \textquotedblleft{}cow\textquotedblright{}, + \emph{rūpa}, \textquotedblleft{}shape\textquotedblright{}.
\end{cousinweb}

\begin{grammarlens}[title={What you've built}]
One more word for this chapter.
\end{grammarlens}

\subsection*{Your turn}
\begin{itemize}
\raggedright
  \item \emph{Say it:} \emph{gôru}
  \item \emph{Say it:} \emph{gôru}, once more
  \item \emph{Say it:} say \emph{gôru}
  \item \emph{Recall:} say the Bengali for a dog, then the Bengali for a cat, then say \emph{gôru} again
\end{itemize}

\begin{tcolorbox}[breakable,colback=teal!4,colframe=teal!35!black,title={Before you move on}]
What does \bn{গরু} mean? (\textquotedblleft{}A cow\textquotedblright{}.) Say it once more.
\end{tcolorbox}
\section[\texorpdfstring{\bn{ছাগল} (chhāgôl)}{ছাগল (chhāgôl)}]{\bn{ছাগল} (chhāgôl) — a goat}

\label{lesson:BN-C69-chhagol}



\begin{quote}
Before the new one: say the Bengali for a cat, then the Bengali for a cow.
\end{quote}

\subsection*{You'll want to know: \bn{ছাগল}}
\textbf{\bn{ছাগল}} — \emph{chhāgôl} — \textquotedblleft{}a goat\textquotedblright{}.

\begin{cousinweb}
\textbf{\bn{ছাগল}} (\emph{chhāgôl}) is Sanskrit \emph{chāgala}, \textquotedblleft{}of a goat; a goat\textquotedblright{}, taken over whole as a learned loan, a tatsama (CDIAL 4963). CDIAL lists the inherited form \emph{cheli}, \textquotedblleft{}goat\textquotedblright{}, under the related \emph{chagala}.
\end{cousinweb}

\begin{grammarlens}[title={What you've built}]
One more word for this chapter.
\end{grammarlens}

\subsection*{Your turn}
\begin{itemize}
\raggedright
  \item \emph{Say it:} \emph{chhāgôl}
  \item \emph{Say it:} \emph{chhāgôl}, once more
  \item \emph{Say it:} say \emph{chhāgôl}
  \item \emph{Recall:} say the Bengali for a cat, then the Bengali for a cow, then say \emph{chhāgôl} again
\end{itemize}

\begin{tcolorbox}[breakable,colback=teal!4,colframe=teal!35!black,title={Before you move on}]
What does \bn{ছাগল} mean? (\textquotedblleft{}A goat\textquotedblright{}.) Say it once more.
\end{tcolorbox}
\section[\texorpdfstring{\bn{ভেড়া} (bheṛā)}{ভেড়া (bheṛā)}]{\bn{ভেড়া} (bheṛā) — a sheep}

\label{lesson:BN-C69-bhera}



\begin{quote}
Before the new one: say the Bengali for a cow, then the Bengali for a goat.
\end{quote}

\subsection*{You'll want to know: \bn{ভেড়া}}
\textbf{\bn{ভেড়া}} — \emph{bheṛā} — \textquotedblleft{}a sheep\textquotedblright{}.

\begin{cousinweb}
\textbf{\bn{ভেড়া}} (\emph{bheṛā}) is a tadbhava, a word inherited from Sanskrit and reshaped by sound change. It comes from \emph{bheḍra}, \textquotedblleft{}ram\textquotedblright{} (CDIAL 9606). CDIAL takes it, following Przyluski, from Austroasiatic.
\end{cousinweb}

\begin{grammarlens}[title={What you've built}]
One more word for this chapter.
\end{grammarlens}

\subsection*{Your turn}
\begin{itemize}
\raggedright
  \item \emph{Say it:} \emph{bheṛā}
  \item \emph{Say it:} \emph{bheṛā}, once more
  \item \emph{Say it:} say \emph{bheṛā}
  \item \emph{Recall:} say the Bengali for a cow, then the Bengali for a goat, then say \emph{bheṛā} again
\end{itemize}

\begin{tcolorbox}[breakable,colback=teal!4,colframe=teal!35!black,title={Before you move on}]
What does \bn{ভেড়া} mean? (\textquotedblleft{}A sheep\textquotedblright{}.) Say it once more.
\end{tcolorbox}
\section[\texorpdfstring{\bn{পাখি} (pākhi)}{পাখি (pākhi)}]{\bn{পাখি} (pākhi) — a bird}

\label{lesson:BN-C69-pakhi}



\begin{quote}
Before the new one: say the Bengali for a goat, then the Bengali for a sheep.
\end{quote}

\subsection*{You'll want to know: \bn{পাখি}}
\textbf{\bn{পাখি}} — \emph{pākhi} — \textquotedblleft{}a bird\textquotedblright{}.

\begin{cousinweb}
\textbf{\bn{পাখি}} (\emph{pākhi}) is a tadbhava, a word inherited from Sanskrit and reshaped by sound change. It comes from \emph{pakṣin}, \textquotedblleft{}winged; bird\textquotedblright{} (CDIAL 7636).
\end{cousinweb}

\begin{grammarlens}[title={What you've built}]
One more word for this chapter.
\end{grammarlens}

\subsection*{Your turn}
\begin{itemize}
\raggedright
  \item \emph{Say it:} \emph{pākhi}
  \item \emph{Say it:} \emph{pākhi}, once more
  \item \emph{Say it:} say \emph{pākhi}
  \item \emph{Recall:} say the Bengali for a goat, then the Bengali for a sheep, then say \emph{pākhi} again
\end{itemize}

\begin{tcolorbox}[breakable,colback=teal!4,colframe=teal!35!black,title={Before you move on}]
What does \bn{পাখি} mean? (\textquotedblleft{}A bird\textquotedblright{}.) Say it once more.
\end{tcolorbox}
\section[\texorpdfstring{\bn{ইঁদুর} (ĩdur)}{ইঁদুর (ĩdur)}]{\bn{ইঁদুর} (ĩdur) — a mouse}

\label{lesson:BN-C69-idur}



\begin{quote}
Before the new one: say the Bengali for a sheep, then the Bengali for a bird.
\end{quote}

\subsection*{You'll want to know: \bn{ইঁদুর}}
\textbf{\bn{ইঁদুর}} — \emph{ĩdur} — \textquotedblleft{}a mouse\textquotedblright{}.

\begin{cousinweb}
\textbf{\bn{ইঁদুর}} (\emph{ĩdur}) is a tadbhava, a word inherited from Sanskrit and reshaped by sound change. It comes from \emph{indūra}, \textquotedblleft{}rat, mouse\textquotedblright{} (CDIAL 2095). CDIAL takes the word, in its other form \emph{undura}, from Austroasiatic.
\end{cousinweb}

\begin{grammarlens}[title={What you've built}]
One more word for this chapter.
\end{grammarlens}

\subsection*{Your turn}
\begin{itemize}
\raggedright
  \item \emph{Say it:} \emph{ĩdur}
  \item \emph{Say it:} \emph{ĩdur}, once more
  \item \emph{Say it:} say \emph{ĩdur}
  \item \emph{Recall:} say the Bengali for a sheep, then the Bengali for a bird, then say \emph{ĩdur} again
\end{itemize}

\begin{tcolorbox}[breakable,colback=teal!4,colframe=teal!35!black,title={Before you move on}]
What does \bn{ইঁদুর} mean? (\textquotedblleft{}A mouse\textquotedblright{}.) Say it once more.
\end{tcolorbox}
